\section{Common Errors}\label{ch5sec:errors}

Even with a working template, errors can appear as you add your own content. This section lists the most common problems and how to resolve them.

\begin{figure}[h]
\centering
\subfloat[][]{\begin{tikzpicture}
\draw[thick, primary, fill=primary!10] (0,0) rectangle (3,1.5);
\node[align=center, font=\small] at (1.5,0.75) {Missing file};
\end{tikzpicture}}
\hfill
\subfloat[][]{\begin{tikzpicture}
\draw[thick, primary, fill=primary!10] (0,0) rectangle (3,1.5);
\node[align=center, font=\small] at (1.5,0.75) {Undefined citation};
\end{tikzpicture}}
\caption{Common error categories.}
\label{fig:build-errors}
\end{figure}

\subsection{Missing Input Files}

If \LaTeX{} reports \texttt{! LaTeX Error: File `...' not found}, check that the path in your \verb|\input| or \verb|\includegraphics| command is correct and that the file exists. Remember that paths are relative to the main \texttt{main.tex} file, not to the chapter file that contains the \verb|\input|.

\subsection{Undefined Citations}

Warnings about undefined citations usually mean a key is missing from the bibliography file. Placeholder keys such as \texttt{placeholder2024demo} are defined in \texttt{class/bibliography/short.bib} as examples; replace them with real entries to remove the warnings.

\begin{table}[h]
\caption{Common error categories and fixes}\label{ch5tab:demo-errors}
\centering
\begin{tabular}{l l}
\toprule
\thead{Error category} & \thead{Fix} \\
\midrule
Missing input file & Verify the path relative to \texttt{main.tex} \\
Undefined citation & Add the entry to the bibliography file \\
Undefined reference & Run \LaTeX{} again or check the label \\
\bottomrule
\end{tabular}
\end{table}

Table~\ref{ch5tab:demo-errors} summarises the most common error categories and their fixes.
